\section{Rechecking the literal source remainder constants}
\label{sec:source-repair}

The original author source remains untouched. Polymath's Proposition
\texttt{RTN-prop} begins at original source line 743. Its proof at
lines 769--776 puts the Stirling error into a numerator $1/3$ where
$2/3$ is required. Subsequent Gaussian displays, including lines
792 and 853--857 (label \texttt{delta1}), replace some occurrences
of $T-3$ by $T'-0.33$ without a valid equality. The passage from
the choice of $\widetilde K_u$ at line 808 to the bounds at lines
824--844 also needs to retain a remainder when
$\widetilde K_u=\lfloor T'/\pi\rfloor$. This section repairs all
three issues and proves that the final literal constants survive.

\begin{theorem}[PMH11: corrected source remainder proof]
\label{PMH11}
For $0\le\sigma\le1$, $T\ge100$, $0<t\le1/2$, put
$T'=T+\pi t/8$, $a=\sqrt{T'/(2\pi)}$, $N=\lfloor a\rfloor$,
and use the original $p,U,C_0(p)$ of
\eqref{PMH1-phase}--\eqref{PMH1-C0}. The original $R_{t,N}$
defined in source \texttt{RTN-def} obeys
\begin{equation}\label{PMH11-final}
 \left|\frac{R_{t,N}(\sigma+iT)}
 {(-1)^{N-1}Ue^{\pi i\sigma/4}e^{t\pi^2/64}M_0(iT')}-C_0(p)\right|
 \le\left(\frac{0.397\,9^\sigma}{a-0.865}
               +\frac5{3(T-6)}\right)e^{3.49/(T-4)}.
\end{equation}
\end{theorem}
\begin{proof}
The imported Riemann--Siegel input is Polymath's Proposition
\texttt{arias-prop}, which attributes the bounds to J.~Arias de
Reyna. Its original labels \texttt{ck-bound-1},
\texttt{ck-bound-2}, \texttt{rsk-bound-1}, and
\texttt{rsk-bound-2} state, for $k\ge1$,
\[
 |C_k(p,u)|\le\frac{\sqrt2}{2\pi}
       \frac{9^u\Gamma(k/2)}{2^k}\quad(u>0),
\]
\[
 |C_k(p,u)|\le\frac{2^{1/2-u}}{2\pi}
       \frac{\Gamma(k/2)}{2\pi((3-2\log2)\pi)^{k/2}}
       \quad(u\le0),
\]
\[
 |RS_K(u+iT')|\le\frac17 2^{3u/2}\Gamma((K+1)/2)(1.1/a)^{K+1}
       \quad(u\ge0),
\]
\[
 |RS_K(u+iT')|\le\frac12(9/10)^{\lceil-u\rceil}
       \Gamma((K+1)/2)(1.1/a)^{K+1}
       \quad(u<0,\ K+u\ge2).
\]
We also retain the source input $|C_0(p)|\le1/2$. All subsequent
integrations and numerical majorants are proved below.

Set $D=T-3$ and $c=11/12$. The source effective Stirling estimate
\texttt{elem-lem}(v) bounds its exponent error at $u+iT'$ by
$1/(6(|u+iT'|-0.66))\le1/(6(T'-0.66))$.
Taylor's formula for \eqref{PMH1-logM} and the source derivative
bound \texttt{alpha-deriv-bound} give
\[
 |\log M_0(u+iT')-\log M_0(iT')-u\alpha(iT')|
 \le\frac{u^2}{4(T'-3)}.
\]
The exact rational terms of $\alpha$ give
\[
 \left|\alpha(iT')-\log a-\frac{\pi i}{4}\right|
 =\left|\frac1{2iT'}+\frac1{iT'-1}\right|\le\frac3{2T'}.
\]
The combined exponent error is therefore at most
\begin{equation}\label{PMH11-pointwise}
 \frac{u^2}{4(T'-3)}+\frac{3|u|}{2T'}+\frac1{6(T'-0.66)}
 \le\frac{u^2+6|u|+2/3}{4D}
 \le\frac{u^2+c}{D}.
\end{equation}
The last step is $6|u|\le3u^2+3$. The numerator $2/3$, not
$1/3$, accounts for the complete Stirling error.

For $u>0$ use $K_u=1$. For $u<0$ use
$K_u=\max(\lfloor-u\rfloor+3,\lfloor2a^2\rfloor)$, so
$K_u+u>2$. Put
\[
 S(u)=\sum_{k=1}^{K_u}\frac{|C_k(p,u)|}{a^k}
            +|RS_{K_u}(u+iT')|.
\]
For $u>0$, the imported bounds give
\[
 S(u)\le\frac{0.2\,9^u}{a}+\frac{0.173\,2^{3u/2}}{a^2}
 \le\frac{0.2\,9^u}{a-0.865}.
\]
Here $2^{3u/2}\le9^u$, $0.2(0.865)=0.173$, and the geometric
series bounds the two reciprocal powers. The value at the single
point $u=0$ has no effect on any integral below.

For $u<0$, the exact constants satisfy
$\sqrt2/(4\pi^2)<0.036$ and
$((3-2\log2)\pi)^{-1/2}<0.445$. Moreover
\[
 \frac{\Gamma((k+2)/2)}{a^{k+2}}
 =\frac{k}{2a^2}\frac{\Gamma(k/2)}{a^k}.
\]
For indices up to $2a^2$ this permits geometric majorants in the
even and odd subsequences. They give
\begin{equation}\label{PMH11-low}
 \sum_{1\le k\le2a^2}\frac{|C_k(p,u)|}{a^k}
 \le0.036\,2^{-u}
 \left(\frac{0.445\sqrt\pi}{a}+\frac{0.247}{a^2}
                         +\frac{0.098}{a^3}\right)
 \le\frac{0.0285\,2^{-u}}{a-0.353}.
\end{equation}
The first constants follow from
\[
 \frac{0.445^2}{1-0.445^2}<0.247,\qquad
 \frac{0.445^3\sqrt\pi}{2(1-0.445^2)}<0.098.
\]
For the last inequality compare the first three coefficients of
the geometric series for $0.0285/(a-0.353)$, using
\[
 0.036(0.445)\sqrt\pi<0.0285,\quad
 0.036(0.247)<0.0285(0.353),\quad
 0.036(0.098)<0.0285(0.353)^2.
\]

The remainder when $K_u=\lfloor2a^2\rfloor$ must still be kept.
Let $r=K_u+1\ge32$; then $a^2\ge(r-1)/2$. Its imported bound,
multiplied by $a2^u$, is at most
\[
 v_r=\frac12(1.1)^r\Gamma(r/2)
       \left(\frac{r-1}{2}\right)^{(1-r)/2}.
\]
Indeed $(9/10)^{\lceil-u\rceil}\le1\le2^{-u}$.
The ratio of successive same-parity terms is exactly
\[
 \frac{v_{r+2}}{v_r}=1.21\frac r{r+1}
       \left(\frac{r-1}{r+1}\right)^{(r-1)/2}
 \le1.21e^{-(r-1)/(r+1)}\le1.21e^{-31/33}<1.
\]
The finite checks $v_{32}<4.897\times10^{-6}$ and
$v_{33}<3.268\times10^{-6}$ are both below $0.0005$.
Thus this complete branch of the remainder is at most
$0.0005\,2^{-u}/a\le0.0005\,2^{-u}/(a-0.353)$.
Together with \eqref{PMH11-low} it is covered by
$0.029\,2^{-u}/(a-0.353)$; no remainder has been dropped.

If $K_u>\lfloor2a^2\rfloor$, the high coefficients and final
remainder all have integer indices in $2a^2\le k\le-u+4$.
The coefficient inequality $0.036(0.445)^k\le(1/2)(1.1)^k$
puts both into one positive sum. Adding unused positive terms in
the two sign ranges gives the uniform bound
\begin{equation}\label{PMH11-S}
 S(u)\le\frac{0.2\,9^u}{a-0.865}+\frac{0.029\,2^{-u}}{a-0.353}
       +\frac{2^{-u}}2
          \sum_{2a^2\le k\le-u+4}(1.1)^k\frac{\Gamma(k/2)}{a^k}.
\end{equation}

Now use the original shift \texttt{RTN-def} with $\beta_N=\pi i/4$
and $u=\sigma+\sqrt t\,v$. The shift keeps $e^{t\pi^2/64}$;
the factors $a^{-u}$ and $e^{u\log a}$ cancel; the phases
$e^{\pi iu/4}$ and $e^{-\pi i\sqrt t v/4}$ leave
$e^{\pi i\sigma/4}$. With
$\mathrm d\nu(u)=(\pi t)^{-1/2}e^{-(u-\sigma)^2/t}\dd u$,
the left side of \eqref{PMH11-final} is bounded by
\begin{equation}\label{PMH11-epsdelta}
 \tfrac12\epsilon+\delta,\quad
 \epsilon=\int_\R(e^{(u^2+c)/D}-1)\mathrm d\nu(u),\quad
 \delta=\int_\R e^{(u^2+c)/D}S(u)\mathrm d\nu(u).
\end{equation}
This uses \eqref{PMH11-pointwise}, $|e^w-1|\le e^{|w|}-1$ and
the unchanged $|C_0|\le1/2$.

For real $l$, completing the square gives the exact identity
\begin{equation}\label{PMH11-Gaussian}
 J(l):=\int_\R e^{(u^2+c)/D+2lu}\mathrm d\nu(u)
 =e^{2l\sigma+tl^2}(1-t/D)^{-1/2}
       \exp\left(\frac cD+\frac{(\sigma+tl)^2}{D-t}\right).
\end{equation}
Every occurrence of $D$ is $T-3$. Expanding the quadratic exponent
gives leading coefficient $-(D-t)/(tD)$; integrating that Gaussian
proves the identity. Since $-\log(1-v)\le v/(1-v)$ for
$0\le v<1$,
\[
 \log J(0)\le\frac{c+\sigma^2+t/2}{D-t}
 \le\frac{13/6}{T-3.5}\le\frac{10}{3(T-6)}.
\]
With $v=10/(3(T-6))$, $e^v-1\le ve^v$, hence
\begin{equation}\label{PMH11-epsilon}
 \tfrac12\epsilon\le\frac5{3(T-6)}e^{3.49/(T-4)}.
\end{equation}
The exponent comparison $10/(3(T-6))\le3.49/(T-4)$ is verified
at $T=100$ after cross multiplication, with positive linear slope
for larger $T$.

The first two terms of \eqref{PMH11-S} integrate to
$\delta_1=0.2J(\log3)/(a-0.865)$ and
$\delta_2=0.029J(-\log\sqrt2)/(a-0.353)$.
Put
\[
 A=\frac{11}{12}+(1+\tfrac12\log3)^2+\frac14
       =3.5670161955\ldots.
\]
After its factor $9^\sigma e^{t\log^23}$, the remaining
logarithm in $J(\log3)$ is at most
$A/(T-3.5)\le A/(T-4)$. The finite strict check
\begin{equation}\label{PMH11-absorb}
 0.2\exp\left(\frac{\log^23}{2}+\frac{A-3.49}{96}\right)
 <0.365986<0.366
\end{equation}
and $T-4\ge96$ prove
\[
 \delta_1\le\frac{0.366\,9^\sigma}{a-0.865}e^{3.49/(T-4)}.
\]
For $M=\log\sqrt2<1$, $|\sigma-tM|\le1$ and
$0.029e^{M^2/2}<0.031$. The remaining exponent in $J(-M)$ is
at most $(13/6)/(T-3.5)<3.49/(T-4)$, so
\[
 \delta_2\le\frac{0.031\,2^{-\sigma}}{a-0.353}e^{3.49/(T-4)}.
\]

The last term of \eqref{PMH11-S}, by Fubini--Tonelli, is exactly
\[
 \delta_3=\frac1{2\sqrt{\pi t}}
 \sum_{k\ge2a^2}\frac{(1.1)^k\Gamma(k/2)}{a^k}
 \int_{-\infty}^{4-k}
 \exp\left(\frac{u^2+c}{D}-\frac{(u-\sigma)^2}{t}
                          -u\log2\right)\dd u.
\]
Here $k\ge32$, hence $k\ge14$ and $u\le-10$. Since
$\sigma\ge0$, $t\le1/2$, $D\ge97$, and
$(1+c/100)/97<1$, this exponent is at most
$-u^2/(2t)-u\log2$. For $u\le4-k$ one has
$-u^2/(2t)\le(k-4)u/(2t)$. Integration and
$((k-4)/(2t)-\log2)\ge(k-6)/(2t)$ give
\[
 \delta_3\le\frac{\sqrt t}{\sqrt\pi}
 \sum_{k\ge2a^2}\frac{(1.1)^k\Gamma(k/2)}{(k-6)a^k}
            e^{-(k-4)^2+(k-4)\log2}.
\]
Set $f_k=(1.1)^k\Gamma(k/2)e^{-(k-4)^2+(k-4)\log2}$.
The two finite bases satisfy $f_{14}<1.042\times10^{-37}$ and
$f_{15}<4.516\times10^{-46}$, both below $10^{-30}$, and
\[
 \frac{f_{k+2}}{f_k}=2.42k e^{12-4k}<1\quad(k\ge14).
\]
Indeed $k e^{-4k}$ decreases there and the case $k=14$ passes.
As $a>3$, enlarging the sum and discarding factors at most one
now gives
\begin{equation}\label{PMH11-delta3}
 \delta_3\le\sum_{k\ge14}\frac{10^{-30}}{a^k}
       \le\frac{2\times10^{-30}}{a^{14}}.
\end{equation}
The complete denominator slack covers this term:
\[
 \frac{0.031\,2^{-\sigma}}{a-0.865}
 -\frac{0.031\,2^{-\sigma}}{a-0.353}
 =\frac{0.031(0.512)2^{-\sigma}}{(a-0.865)(a-0.353)}
 \ge\frac{0.007936}{a^2}
 >\frac{2\times10^{-30}}{a^{14}}.
\]
Thus
\[
 \delta\le\frac{0.366\,9^\sigma+0.031\,2^{-\sigma}}
                  {a-0.865}e^{3.49/(T-4)}
 \le\frac{0.397\,9^\sigma}{a-0.865}e^{3.49/(T-4)}.
\]
Together with \eqref{PMH11-epsdelta} and \eqref{PMH11-epsilon}
this proves \eqref{PMH11-final}. Every finite transcendental
comparison in this proof is checked in
\texttt{corrected\_rtn\_constant\_checks()}.
\end{proof}
